\documentclass[11pt]{article}
\usepackage[margin=1in]{geometry}
\usepackage{amsmath,amssymb,amsthm,mathtools}
\usepackage{enumitem}
\usepackage{hyperref}
\hypersetup{colorlinks=true,linkcolor=blue,citecolor=blue,urlcolor=blue}

\newtheorem{theorem}{Theorem}
\newtheorem{proposition}{Proposition}
\newtheorem{lemma}{Lemma}
\newtheorem{corollary}{Corollary}
\newtheorem{definition}{Definition}
\newtheorem{hypothesis}{Hypothesis}
\newtheorem{remark}{Remark}
\newtheorem{warning}{Warning}
\newtheorem{conjecture}{Conjecture}
\newtheorem{problem}{Problem}

\newcommand{\Z}{\mathbb Z}
\newcommand{\Q}{\mathbb Q}
\newcommand{\F}{\mathbb F}
\newcommand{\C}{\mathbb C}
\newcommand{\G}{\Gamma}
\newcommand{\E}{\mathbb E}
\newcommand{\1}{\mathbf 1}
\newcommand{\ord}{\operatorname{ord}}
\newcommand{\lcm}{\operatorname{lcm}}
\newcommand{\im}{\operatorname{im}}
\newcommand{\rank}{\operatorname{rank}}
\newcommand{\Gal}{\operatorname{Gal}}

\title{On Some Tempting but Insufficient Routes to Fixed-Radical Kummer Non-Concentration}
\author{Jared Wilder}
\date{Draft 0.1 --- 2026-05-12}

\begin{document}
\maketitle

\begin{abstract}
We record several plausible but insufficient routes to a fixed-radical Kummer non-concentration theorem arising in a conditional approach to Erd\H{o}s--Graham Problem \#203.  The purpose is defensive and clarifying: to separate correct structural lemmas from missing analytic distribution estimates.  The note identifies the high-\(\ell\) prime Hecke non-concentration seam as the precise remaining analytic problem.
\end{abstract}

\section{The target theorem}

For \(\ell\le(\log Q)^B\) and \((r,s)\ne(0,0)\in\F_\ell^2\), define
\[
\theta=2^r3^s.
\]
The desired non-concentration bound is
\[
\#\{q\le Q:q\equiv1\bmod\ell,\ \theta\in(\F_q^*)^\ell\}
\le
(1-\ell^{-C})\pi(Q;\ell,1).
\]
This is PHNC.

\section{Correct structural facts}

\subsection{Kummer independence}

For \((r,s)\ne0\), the element \(2^r3^s\) is not an \(\ell\)-th power in \(\Q(\zeta_\ell)\).  This follows from valuations at primes above \(2\) and \(3\).  Therefore the Kummer layer
\[
\Q(\zeta_\ell,(2^r3^s)^{1/\ell})/\Q(\zeta_\ell)
\]
has degree \(\ell\).

\subsection{Order and non-real character}

The corresponding Kummer character has order exactly \(\ell\).  If \(\ell\ge3\), it is non-real.

\subsection{Power-stability of small deficit}

Let \(f:X\to S^1\) and
\[
\delta(f)=1-|\E f|.
\]
Then
\[
\delta(f^m)\le m^2\delta(f).
\]
Consequently, if an order-\(\ell\) character has deficit \(\le\ell^{-C}\), \(C\ge3\), then its entire nontrivial cyclic line has deficit \(\le1/2\) for large \(\ell\).

\section{Tempting route 1: BFI as a bad-set bound}

A tempting claim is that Bombieri--Friedlander--Iwaniec type results imply a bound of the shape
\[
|B_{1/2}|\ll \frac{\log Q}{\ell}
\]
for the set of Kummer characters with deficit at most \(1/2\).

\begin{warning}
This is not a valid plug-in citation.  The BFI theorem concerns primes in arithmetic progressions with well-factorable weights.  It does not directly supply this fixed-radical Hecke/Kummer bad-character bound uniformly in \(\ell\).
\end{warning}

\subsection{What the bound would imply if true}

\begin{proposition}
Suppose one had
\[
|B_{1/2}|\le C_0\frac{\log Q}{\ell}.
\]
Then in the regime
\[
\ell^2>C_0\log Q\cdot\frac{\ell}{\ell-1},
\]
PHNC would follow.
\end{proposition}

\begin{proof}
If PHNC failed, there would be a character \(\chi\) with \(\delta(\chi)\le\ell^{-C}\).  The power-stability lemma implies that the \(\ell-1\) nontrivial powers \(\chi,\dots,\chi^{\ell-1}\) all lie in \(B_{1/2}\).  Hence \(|B_{1/2}|\ge\ell-1\).  The assumed bound gives \(|B_{1/2}|<\ell-1\), contradiction.
\end{proof}

The structural part is correct.  The missing part is the analytic bad-set bound.

\section{Tempting route 2: Pl\"unnecke--Ruzsa}

Another tempting route is to use additive combinatorics on the set of bad characters.

\begin{warning}
Pl\"unnecke--Ruzsa does not by itself provide prime distribution.  In this setting, the one-character cyclic-line implication is already obtained by the elementary \(L^2\) power-stability lemma.  The hard part is still the analytic bound on how many medium-deficit characters exist.
\end{warning}

\section{Tempting route 3: fixed-field ineffective Chebotarev}

For each fixed \(\ell,r,s\), Chebotarev gives eventual equidistribution:
\[
\#\{q\le Q:q\equiv1\bmod\ell,\ \theta\in(\F_q^*)^\ell\}
\sim
\frac1\ell\pi(Q;\ell,1).
\]

\begin{warning}
This does not imply PHNC uniformly for the growing family
\[
\ell\le(\log Q)^B.
\]
An ineffective threshold depending on the individual field does not automatically yield a single \(Q_0(B,C)\) valid for all \(\ell\le(\log Q)^B\).
\end{warning}

\section{Tempting route 4: multiplicative-order lower bounds}

One might try to use lower bounds on \(\ord_q(2^r3^s)\).

\begin{warning}
The condition that \(2^r3^s\) is an \(\ell\)-th power modulo \(q\) is equivalent to
\[
\ord_q(2^r3^s)\mid\frac{q-1}{\ell}.
\]
For polylogarithmic \(\ell\), the upper bound \((q-1)/\ell\) is still very large.  General lower bounds for multiplicative order do not rule out this divisibility for almost all \(q\equiv1\bmod\ell\).
\end{warning}

\section{Tempting route 5: generic Chebotarev--Brun--Titchmarsh}

Chebotarev--Brun--Titchmarsh type estimates are often useful for upper bounds on splitting primes.

\begin{warning}
Known forms carry field-invariant dependencies.  In the present growing cyclotomic/Kummer family, those dependencies again restrict the usable range.  They do not automatically prove the full high-\(\ell\) seam
\[
(\log Q)^{1-\varepsilon}<\ell\le(\log Q)^B.
\]
\end{warning}

\section{The actual open seam}

The remaining analytic problem is:

\begin{problem}[High-seam fixed-radical Kummer non-concentration]
For every fixed \(B>0\), prove that there exists \(C=C(B)>0\) such that, uniformly for all primes
\[
\ell\le(\log Q)^B
\]
and all nonzero \((r,s)\in\F_\ell^2\),
\[
\#\{q\le Q:q\equiv1\bmod\ell,\ 2^r3^s\in(\F_q^*)^\ell\}
\le
(1-\ell^{-C})\pi(Q;\ell,1).
\]
\end{problem}

\section{Conclusion}

The failed routes are useful because they isolate the exact gap.  The structural algebra is real; the high-\(\ell\) analytic distribution theorem remains open.

\bibliographystyle{alpha}
\begin{thebibliography}{9}

\bibitem{BFI}
E. Bombieri, J. Friedlander, and H. Iwaniec.
\newblock Primes in arithmetic progressions to large moduli.
\newblock \emph{Acta Mathematica}, 156:203--251, 1986.

\bibitem{LO}
J. C. Lagarias and A. M. Odlyzko.
\newblock Effective versions of the Chebotarev density theorem.
\newblock \emph{Algebraic Number Fields}, 1977.

\bibitem{TZ}
J. Thorner and A. Zaman.
\newblock A unified and improved Chebotarev density theorem.
\newblock \emph{Algebra \& Number Theory}, 13(5):1039--1068, 2019.

\end{thebibliography}

\end{document}
